\documentclass[11pt,a4paper]{article}
\usepackage[utf8]{inputenc}
\usepackage[T1]{fontenc}
\usepackage[spanish,es-minimal,es-noshorthands]{babel}
\AtBeginDocument{\renewcommand{\sin}{\operatorname{sen}}\renewcommand{\arcsin}{\operatorname{arcsen}}}
\usepackage{lmodern}
\usepackage[sfdefault,lining]{FiraSans}
\usepackage[a4paper,top=2.3cm,bottom=2.2cm,left=1.8cm,right=1.8cm,headheight=1.3cm,headsep=0.4cm,footskip=1cm]{geometry}
\usepackage{amsmath,amssymb}
\usepackage{icomma}
\usepackage{graphicx}
\usepackage[table]{xcolor}
\usepackage{tikz}
\usepackage{fontawesome5}
\usepackage{ccicons}
\usepackage{fancyhdr}
\usepackage{lastpage}
\usepackage{tasks}
\usepackage{needspace}
\usepackage[most]{tcolorbox}
\usepackage{array}
\usepackage[hidelinks]{hyperref}
\AtBeginDocument{\hypersetup{pdftitle={Soluciones de la ficha de la unidad 8: Áreas y perímetros. 1º ESO},pdfsubject={Matemáticas 1º ESO},pdfkeywords={perímetro, área, unidades de longitud, unidades de superficie, hectárea, cuadrado, rectángulo, romboide, triángulo, rombo, trapecio, polígonos regulares, apotema, circunferencia, número pi, círculo, sector circular, corona circular, figuras compuestas, teorema de Pitágoras, catetos, hipotenusa}}}

\definecolor{navy}{HTML}{1B2748}
\definecolor{accent}{HTML}{E8603A}
\definecolor{teal}{HTML}{1F9E90}
\definecolor{soft}{HTML}{EEF1F8}
\definecolor{muted}{HTML}{5D6685}

\setlength{\parindent}{0pt}
\setlength{\parskip}{3pt}
\settasks{label=\alph*), label-format=\bfseries\color{navy}, label-width=1.3em, item-indent=1.8em, column-sep=0.6em, after-item-skip=5pt, before-skip=3pt, after-skip=4pt}

% ---------- capçalera i peu ----------
\newcommand{\llicencia}{\href{https://creativecommons.org/licenses/by-nc-sa/4.0/deed.es}{\ccbyncsa}\ \ Licencia Creative Commons Atribución-NoComercial-CompartirIgual 4.0 Internacional (CC BY-NC-SA 4.0)}
\pagestyle{fancy}
\fancyhf{}
\fancyhead[L]{\raisebox{-0.35\height}{\includegraphics[height=0.95cm]{logo-gva.png}}}
\fancyhead[C]{\small\color{muted}Soluciones de la ficha · Unidad 8 · 1\textsuperscript{o} ESO}
\fancyhead[R]{\raisebox{-0.35\height}{\includegraphics[height=0.95cm]{logo-vermellar.png}}}
\fancyfoot[L]{\scriptsize\color{muted}\llicencia}
\fancyfoot[R]{\small\color{muted}\thepage}
\renewcommand{\headrulewidth}{0.4pt}
\renewcommand{\headrule}{\color{muted!40}\hrule height \headrulewidth}
\fancypagestyle{primera}{\fancyhf{}\renewcommand{\headrulewidth}{0pt}
  \fancyfoot[L]{\scriptsize\color{muted}\llicencia}
  \fancyfoot[R]{\small\color{muted}\thepage}}

% ---------- nivells i exercicis ----------
\newcommand{\estrelles}[1]{{\color{accent}\ifcase#1\or\faStar\faStar[regular]\faStar[regular]\or\faStar\faStar\faStar[regular]\or\faStar\faStar\faStar\fi}}
\newcounter{exer}
\newcommand{\ex}[1]{\par\needspace{4\baselineskip}\medskip\stepcounter{exer}\noindent{\color{navy}\bfseries\theexer.}\ \ {\small\estrelles{#1}}\ \ }

\newcommand{\seccio}[2]{\par\needspace{10\baselineskip}\bigskip
  \begin{tcolorbox}[colback=navy,colframe=navy,coltext=white,arc=2mm,boxrule=0pt,left=3mm,right=3mm,top=1.2mm,bottom=1.2mm]
  {\large\bfseries #1\quad #2}
  \end{tcolorbox}\vspace{-2pt}}

\newcolumntype{C}[1]{>{\centering\arraybackslash}p{#1}}

% quadre de recordatori (només a la fitxa d'exercicis)
\newcommand{\recorda}[1]{\par\smallskip\begin{tcolorbox}[colback=accent!7,colframe=accent!55,arc=2mm,boxrule=0.5pt,left=2.5mm,right=2.5mm,top=1mm,bottom=1mm]\small{\color{accent}\faLightbulb[regular]}\ \textbf{Recuerda:}\ #1\end{tcolorbox}}
\newcommand{\res}[1]{\boxed{#1}}\newcommand{\rest}[1]{\fbox{#1}}

\def\fmneed{14}
% =====================================================================
% Capçalera FentMates de les fitxes: un color per a cada curs.
% Els colors són els de la pàgina inicial (LEVELS a assets/js/materials.js);
% el logo ve de _eines/logo/pdf/ (python3 _eines/logo/genera.py).
%
% Ús, al final del preàmbul de la unitat:
%   % =====================================================================
% Capçalera FentMates de les fitxes: un color per a cada curs.
% Els colors són els de la pàgina inicial (LEVELS a assets/js/materials.js);
% el logo ve de _eines/logo/pdf/ (python3 _eines/logo/genera.py).
%
% Ús, al final del preàmbul de la unitat:
%   % =====================================================================
% Capçalera FentMates de les fitxes: un color per a cada curs.
% Els colors són els de la pàgina inicial (LEVELS a assets/js/materials.js);
% el logo ve de _eines/logo/pdf/ (python3 _eines/logo/genera.py).
%
% Ús, al final del preàmbul de la unitat:
%   \input{../comu/fentmates.tex}
%   \fmcurs{1eso}
% i, al començament del document:
%   \capcalera{Matemàtiques}{1r ESO}{Fitxa d'exercicis · Unitat 2}{Divisibilitat}
% Opcional abans de \input: \def\fmneed{14} (espai mínim abans de cada secció).
% =====================================================================
\makeatletter
\def\fm@logodir{../../logo/pdf}
\providecommand{\fmneed}{10}
\def\fm@cursnom{}

% clau del curs, color i color del text damunt del color
\newcommand{\fm@curs}[3]{\@namedef{fm@color@#1}{#2}\@namedef{fm@text@#1}{#3}}
\fm@curs{1eso}{E6B84A}{navy}
\fm@curs{2eso}{3B82F6}{white}
\fm@curs{3eso}{FF7A50}{white}
\fm@curs{4eso-a}{A78BFA}{white}
\fm@curs{4eso-b}{6D28D9}{white}
\fm@curs{1batx-mi}{9B1C1C}{white}
\fm@curs{1batx-ccss}{EF6B6B}{white}
\fm@curs{2batx-mii}{1E6B47}{white}
\fm@curs{2batx-ccss}{2BB3A3}{white}

\newcommand{\fmcurs}[1]{%
  \@ifundefined{fm@color@#1}{\PackageError{fentmates}{Curs desconegut: #1}{}}{}%
  \edef\fm@tmp{\noexpand\definecolor{curs}{HTML}{\@nameuse{fm@color@#1}}%
               \noexpand\colorlet{curstext}{\@nameuse{fm@text@#1}}}\fm@tmp
  \colorlet{cursfosc}{curs!72!black}%
  \def\fm@logo{\fm@logodir/vertical-#1.pdf}%
  \def\fm@marca{\fm@logodir/marca-#1.pdf}}

% ---------- capçalera de la primera pàgina ----------
% #1 matèria (petit), #2 curs, #3 tipus de document i unitat, #4 títol
\newcommand{\capcalera}[4]{%
  \gdef\fm@cursnom{#2}%
  \par\noindent
  \begin{tcolorbox}[colback=curs,colframe=curs,boxrule=0pt,arc=2.2mm,width=3.5cm,height=2.15cm,
      valign=center,left=3mm,right=2mm,top=1mm,bottom=1mm,nobeforeafter,box align=center]
    \color{curstext}{\footnotesize #1}\par\vspace{1pt}{\fontsize{19}{21}\selectfont\bfseries #2\par}%
  \end{tcolorbox}\hspace{5mm}%
  \begin{minipage}[c]{\dimexpr\linewidth-3.5cm-5mm-2.9cm\relax}
    {\color{cursfosc}\small\bfseries\MakeUppercase{#3}}\par\vspace{2pt}
    {\color{navy}\fontsize{20}{23}\selectfont\bfseries #4\par}
  \end{minipage}\hfill
  \begin{minipage}[c]{2.6cm}\raggedleft\includegraphics[width=2.6cm]{\fm@logo}\end{minipage}\par
  \vspace{2pt}}

% ---------- barres de secció del color del curs ----------
\renewcommand{\seccio}[2]{\par\needspace{\fmneed\baselineskip}\bigskip
  \begin{tcolorbox}[colback=curs,colframe=curs,coltext=curstext,arc=2mm,boxrule=0pt,left=3mm,right=3mm,top=1.2mm,bottom=1.2mm]
  {\large\bfseries #1\quad #2}
  \end{tcolorbox}\vspace{-2pt}}

% ---------- estrelles de dificultat del color del curs (a Batxillerat no se'n posen) ----------
\@ifundefined{estrelles}{}{%
  \renewcommand{\estrelles}[1]{{\color{curs}\ifcase#1\or\faStar\faStar[regular]\faStar[regular]\or\faStar\faStar\faStar[regular]\or\faStar\faStar\faStar\fi}}}

% ---------- capçalera i peu de les pàgines ----------
% primera pàgina: logos institucionals al peu
\fancypagestyle{primera}{\fancyhf{}\renewcommand{\headrulewidth}{0pt}
  \fancyfoot[L]{\raisebox{-0.32\height}{\includegraphics[height=0.78cm]{logo-gva.png}}\hspace{5mm}%
    \raisebox{-0.32\height}{\includegraphics[height=0.78cm]{logo-vermellar.png}}}
  \fancyfoot[R]{\scriptsize\color{muted}\href{https://creativecommons.org/licenses/by-nc-sa/4.0/deed.ca}{\ccbyncsa}\ \ CC BY-NC-SA 4.0 · fentmates.org\qquad\small\thepage}}
% la resta: el símbol a l'esquerra i el curs a la dreta, amb el títol al centre com abans
\fancyhead[L]{\raisebox{-0.3\height}{\includegraphics[height=0.62cm]{\fm@marca}}}
\fancyhead[R]{\small\bfseries\color{cursfosc}\fm@cursnom}
\renewcommand{\headrule}{\color{curs!45}\hrule height \headrulewidth}
\makeatother

%   \fmcurs{1eso}
% i, al començament del document:
%   \capcalera{Matemàtiques}{1r ESO}{Fitxa d'exercicis · Unitat 2}{Divisibilitat}
% Opcional abans de \input: \def\fmneed{14} (espai mínim abans de cada secció).
% =====================================================================
\makeatletter
\def\fm@logodir{../../logo/pdf}
\providecommand{\fmneed}{10}
\def\fm@cursnom{}

% clau del curs, color i color del text damunt del color
\newcommand{\fm@curs}[3]{\@namedef{fm@color@#1}{#2}\@namedef{fm@text@#1}{#3}}
\fm@curs{1eso}{E6B84A}{navy}
\fm@curs{2eso}{3B82F6}{white}
\fm@curs{3eso}{FF7A50}{white}
\fm@curs{4eso-a}{A78BFA}{white}
\fm@curs{4eso-b}{6D28D9}{white}
\fm@curs{1batx-mi}{9B1C1C}{white}
\fm@curs{1batx-ccss}{EF6B6B}{white}
\fm@curs{2batx-mii}{1E6B47}{white}
\fm@curs{2batx-ccss}{2BB3A3}{white}

\newcommand{\fmcurs}[1]{%
  \@ifundefined{fm@color@#1}{\PackageError{fentmates}{Curs desconegut: #1}{}}{}%
  \edef\fm@tmp{\noexpand\definecolor{curs}{HTML}{\@nameuse{fm@color@#1}}%
               \noexpand\colorlet{curstext}{\@nameuse{fm@text@#1}}}\fm@tmp
  \colorlet{cursfosc}{curs!72!black}%
  \def\fm@logo{\fm@logodir/vertical-#1.pdf}%
  \def\fm@marca{\fm@logodir/marca-#1.pdf}}

% ---------- capçalera de la primera pàgina ----------
% #1 matèria (petit), #2 curs, #3 tipus de document i unitat, #4 títol
\newcommand{\capcalera}[4]{%
  \gdef\fm@cursnom{#2}%
  \par\noindent
  \begin{tcolorbox}[colback=curs,colframe=curs,boxrule=0pt,arc=2.2mm,width=3.5cm,height=2.15cm,
      valign=center,left=3mm,right=2mm,top=1mm,bottom=1mm,nobeforeafter,box align=center]
    \color{curstext}{\footnotesize #1}\par\vspace{1pt}{\fontsize{19}{21}\selectfont\bfseries #2\par}%
  \end{tcolorbox}\hspace{5mm}%
  \begin{minipage}[c]{\dimexpr\linewidth-3.5cm-5mm-2.9cm\relax}
    {\color{cursfosc}\small\bfseries\MakeUppercase{#3}}\par\vspace{2pt}
    {\color{navy}\fontsize{20}{23}\selectfont\bfseries #4\par}
  \end{minipage}\hfill
  \begin{minipage}[c]{2.6cm}\raggedleft\includegraphics[width=2.6cm]{\fm@logo}\end{minipage}\par
  \vspace{2pt}}

% ---------- barres de secció del color del curs ----------
\renewcommand{\seccio}[2]{\par\needspace{\fmneed\baselineskip}\bigskip
  \begin{tcolorbox}[colback=curs,colframe=curs,coltext=curstext,arc=2mm,boxrule=0pt,left=3mm,right=3mm,top=1.2mm,bottom=1.2mm]
  {\large\bfseries #1\quad #2}
  \end{tcolorbox}\vspace{-2pt}}

% ---------- estrelles de dificultat del color del curs (a Batxillerat no se'n posen) ----------
\@ifundefined{estrelles}{}{%
  \renewcommand{\estrelles}[1]{{\color{curs}\ifcase#1\or\faStar\faStar[regular]\faStar[regular]\or\faStar\faStar\faStar[regular]\or\faStar\faStar\faStar\fi}}}

% ---------- capçalera i peu de les pàgines ----------
% primera pàgina: logos institucionals al peu
\fancypagestyle{primera}{\fancyhf{}\renewcommand{\headrulewidth}{0pt}
  \fancyfoot[L]{\raisebox{-0.32\height}{\includegraphics[height=0.78cm]{logo-gva.png}}\hspace{5mm}%
    \raisebox{-0.32\height}{\includegraphics[height=0.78cm]{logo-vermellar.png}}}
  \fancyfoot[R]{\scriptsize\color{muted}\href{https://creativecommons.org/licenses/by-nc-sa/4.0/deed.ca}{\ccbyncsa}\ \ CC BY-NC-SA 4.0 · fentmates.org\qquad\small\thepage}}
% la resta: el símbol a l'esquerra i el curs a la dreta, amb el títol al centre com abans
\fancyhead[L]{\raisebox{-0.3\height}{\includegraphics[height=0.62cm]{\fm@marca}}}
\fancyhead[R]{\small\bfseries\color{cursfosc}\fm@cursnom}
\renewcommand{\headrule}{\color{curs!45}\hrule height \headrulewidth}
\makeatother

%   \fmcurs{1eso}
% i, al començament del document:
%   \capcalera{Matemàtiques}{1r ESO}{Fitxa d'exercicis · Unitat 2}{Divisibilitat}
% Opcional abans de \input: \def\fmneed{14} (espai mínim abans de cada secció).
% =====================================================================
\makeatletter
\def\fm@logodir{../../logo/pdf}
\providecommand{\fmneed}{10}
\def\fm@cursnom{}

% clau del curs, color i color del text damunt del color
\newcommand{\fm@curs}[3]{\@namedef{fm@color@#1}{#2}\@namedef{fm@text@#1}{#3}}
\fm@curs{1eso}{E6B84A}{navy}
\fm@curs{2eso}{3B82F6}{white}
\fm@curs{3eso}{FF7A50}{white}
\fm@curs{4eso-a}{A78BFA}{white}
\fm@curs{4eso-b}{6D28D9}{white}
\fm@curs{1batx-mi}{9B1C1C}{white}
\fm@curs{1batx-ccss}{EF6B6B}{white}
\fm@curs{2batx-mii}{1E6B47}{white}
\fm@curs{2batx-ccss}{2BB3A3}{white}

\newcommand{\fmcurs}[1]{%
  \@ifundefined{fm@color@#1}{\PackageError{fentmates}{Curs desconegut: #1}{}}{}%
  \edef\fm@tmp{\noexpand\definecolor{curs}{HTML}{\@nameuse{fm@color@#1}}%
               \noexpand\colorlet{curstext}{\@nameuse{fm@text@#1}}}\fm@tmp
  \colorlet{cursfosc}{curs!72!black}%
  \def\fm@logo{\fm@logodir/vertical-#1.pdf}%
  \def\fm@marca{\fm@logodir/marca-#1.pdf}}

% ---------- capçalera de la primera pàgina ----------
% #1 matèria (petit), #2 curs, #3 tipus de document i unitat, #4 títol
\newcommand{\capcalera}[4]{%
  \gdef\fm@cursnom{#2}%
  \par\noindent
  \begin{tcolorbox}[colback=curs,colframe=curs,boxrule=0pt,arc=2.2mm,width=3.5cm,height=2.15cm,
      valign=center,left=3mm,right=2mm,top=1mm,bottom=1mm,nobeforeafter,box align=center]
    \color{curstext}{\footnotesize #1}\par\vspace{1pt}{\fontsize{19}{21}\selectfont\bfseries #2\par}%
  \end{tcolorbox}\hspace{5mm}%
  \begin{minipage}[c]{\dimexpr\linewidth-3.5cm-5mm-2.9cm\relax}
    {\color{cursfosc}\small\bfseries\MakeUppercase{#3}}\par\vspace{2pt}
    {\color{navy}\fontsize{20}{23}\selectfont\bfseries #4\par}
  \end{minipage}\hfill
  \begin{minipage}[c]{2.6cm}\raggedleft\includegraphics[width=2.6cm]{\fm@logo}\end{minipage}\par
  \vspace{2pt}}

% ---------- barres de secció del color del curs ----------
\renewcommand{\seccio}[2]{\par\needspace{\fmneed\baselineskip}\bigskip
  \begin{tcolorbox}[colback=curs,colframe=curs,coltext=curstext,arc=2mm,boxrule=0pt,left=3mm,right=3mm,top=1.2mm,bottom=1.2mm]
  {\large\bfseries #1\quad #2}
  \end{tcolorbox}\vspace{-2pt}}

% ---------- estrelles de dificultat del color del curs (a Batxillerat no se'n posen) ----------
\@ifundefined{estrelles}{}{%
  \renewcommand{\estrelles}[1]{{\color{curs}\ifcase#1\or\faStar\faStar[regular]\faStar[regular]\or\faStar\faStar\faStar[regular]\or\faStar\faStar\faStar\fi}}}

% ---------- capçalera i peu de les pàgines ----------
% primera pàgina: logos institucionals al peu
\fancypagestyle{primera}{\fancyhf{}\renewcommand{\headrulewidth}{0pt}
  \fancyfoot[L]{\raisebox{-0.32\height}{\includegraphics[height=0.78cm]{logo-gva.png}}\hspace{5mm}%
    \raisebox{-0.32\height}{\includegraphics[height=0.78cm]{logo-vermellar.png}}}
  \fancyfoot[R]{\scriptsize\color{muted}\href{https://creativecommons.org/licenses/by-nc-sa/4.0/deed.ca}{\ccbyncsa}\ \ CC BY-NC-SA 4.0 · fentmates.org\qquad\small\thepage}}
% la resta: el símbol a l'esquerra i el curs a la dreta, amb el títol al centre com abans
\fancyhead[L]{\raisebox{-0.3\height}{\includegraphics[height=0.62cm]{\fm@marca}}}
\fancyhead[R]{\small\bfseries\color{cursfosc}\fm@cursnom}
\renewcommand{\headrule}{\color{curs!45}\hrule height \headrulewidth}
\makeatother

\fmcurs{1eso}

\begin{document}
\thispagestyle{primera}

\capcalera{Matemáticas}{1\textsuperscript{o} ESO}{Soluciones de la ficha · Unidad 8}{Áreas y perímetros}
\vspace{8pt}

\begin{tcolorbox}[colback=soft,colframe=soft,arc=2mm,left=3mm,right=3mm,top=1.5mm,bottom=1.5mm]
\small Intenta resolver cada ejercicio antes de mirar su solución. Si no te sale el mismo resultado, repasa los pasos o consulta los apuntes interactivos de la unidad: encontrarás herramientas para cambiar de unidades y para explorar las áreas y los perímetros de todas las figuras.
\end{tcolorbox}

% =====================================================
\seccio{1}{Perímetro y área. Unidades de medida}

\ex{1} \begin{tasks}(3)
  \task $A=\res{6\text{ cm}^2}$, $P=\res{10\text{ cm}}$
  \task $A=\res{7\text{ cm}^2}$, $P=\res{16\text{ cm}}$
  \task $A=\res{5\text{ cm}^2}$, $P=\res{12\text{ cm}}$
\end{tasks}

\ex{1} \begin{tasks}(5)
  \task $\res{4700\text{ m}}$
  \task $\res{3,8\text{ m}}$
  \task $\res{2,5\text{ m}}$
  \task $\res{900\text{ m}}$
  \task $\res{1,25\text{ m}}$
\end{tasks}

\ex{1} \begin{tasks}(4)
  \task $\res{320\text{ cm}}$
  \task $\res{4,5\text{ cm}}$
  \task $\res{80\text{ dam}}$
  \task $\res{0,76\text{ hm}}$
\end{tasks}

\ex{1} \begin{tasks}(3)
  \task $\res{500\text{ cm}^2}$
  \task $\res{240\text{ dm}^2}$
  \task $\res{3,5\text{ m}^2}$
  \task $\res{7000\text{ m}^2}$
  \task $\res{800\text{ hm}^2}$
\end{tasks}

\ex{2} \begin{tasks}(4)
  \task $\res{35\,000\text{ m}^2}$
  \task $\res{0,45\text{ ha}}$
  \task $\res{1200\text{ m}^2}$
  \task $\res{25\text{ ha}}$
\end{tasks}

\ex{2} Lo pasamos todo a m$^2$: $50\,000$, $40\,000$, $52\,000$ y $48\,000$.\par\rest{$4\text{ ha}<48\,000\text{ m}^2<0,05\text{ km}^2<520\text{ dam}^2$}

\ex{2} $2,4\text{ ha}=24\,000\text{ m}^2$; \ $24\,000\cdot18=\res{432\,000\text{ €}}$

\ex{3} La base y la altura suman 10: \ $1\times9$ ($A=9$), $2\times8$ ($A=16$), $3\times7$ ($A=21$), $4\times6$ ($A=24$) y $5\times5$ ($A=25$). El área más grande es la del \rest{cuadrado de 5\,cm}: $25\text{ cm}^2$.


% =====================================================
\seccio{2}{Perímetro de los polígonos}

\ex{1} \begin{tasks}(2)
  \task $\res{27,5\text{ cm}}$
  \task $4\cdot8,5=\res{34\text{ cm}}$
  \task $2\cdot15+2\cdot9=\res{48\text{ m}}$
  \task $4\cdot6,3=\res{25,2\text{ cm}}$
\end{tasks}

\ex{1} \begin{tasks}(2)
  \task $5\cdot7=\res{35\text{ cm}}$
  \task $6\cdot4,5=\res{27\text{ cm}}$
  \task $8\cdot3,2=\res{25,6\text{ cm}}$
  \task $10\cdot12=\res{120\text{ mm}}$
\end{tasks}

\ex{1} $P=5+3+4+3,5+3=\res{18,5\text{ cm}}$

\ex{2} \begin{tasks}(3)
  \task $46:4=\res{11,5\text{ cm}}$
  \task $51:6=\res{8,5\text{ cm}}$
  \task $13,5:3=\res{4,5\text{ m}}$
\end{tasks}

\ex{2} Las dos bases miden $40$\,cm; quedan $64-40=24$\,cm para las dos alturas: \ $h=24:2=\res{12\text{ cm}}$.

\ex{2} $P=2\cdot120+2\cdot85=410$\,m. \ Con tres hilos: $3\cdot410=\res{1230\text{ m}}$. \ Precio: $1230\cdot0,35=\res{430,50\text{ €}}$.

\ex{3} El perímetro equivale a $3+1+3+1=8$ alturas: $h=96:8=\res{12\text{ cm}}$ y $b=3\cdot12=\res{36\text{ cm}}$.

\ex{3} El perímetro equivale a $1+2+2=5$ bases: base $=40:5=\res{8\text{ cm}}$ y los lados iguales, \res{16\text{ cm}} cada uno.


% =====================================================
\seccio{3}{Área del cuadrado y del rectángulo}

\ex{1} \begin{tasks}(2)
  \task $A=\res{144\text{ cm}^2}$, $P=\res{48\text{ cm}}$
  \task $A=\res{120\text{ cm}^2}$, $P=\res{46\text{ cm}}$
  \task $A=\res{30\text{ m}^2}$, $P=\res{23\text{ m}}$
  \task $A=\res{6,25\text{ m}^2}$, $P=\res{10\text{ m}}$
\end{tasks}

\ex{1} $h=84:12=\res{7\text{ cm}}$; \ $P=2\cdot12+2\cdot7=\res{38\text{ cm}}$

\ex{1} \begin{tasks}(1)
  \task $11^2=121\Rightarrow l=\res{11\text{ m}}$; \ $P=\res{44\text{ m}}$
  \task $l=28:4=7$\,cm; \ $A=7^2=\res{49\text{ cm}^2}$
\end{tasks}

\ex{2} $A=5,4\cdot3,5=18,9\text{ m}^2$; \ $18,9\cdot24=\res{453,60\text{ €}}$

\ex{2} $6\cdot2,5-1,5\cdot1,2=15-1,8=13,2\text{ m}^2$. \ $13,2:5=2,64$: hay que comprar \rest{3 botes}.

\ex{2} Caben $800:50=16$ baldosas a lo largo y $600:50=12$ a lo ancho: \ $16\cdot12=\res{192}$ baldosas.

\ex{3} $P$: de 20\,cm a 40\,cm, \rest{se duplica}. \ $A$: de 25\,cm$^2$ a 100\,cm$^2$, \rest{se multiplica por 4}.

\ex{3} $1\times36$ ($P=74$), $2\times18$ ($P=40$), $3\times12$ ($P=30$), $4\times9$ ($P=26$), $6\times6$ ($P=24$). El de perímetro más pequeño es el \rest{cuadrado de 6\,cm}.


% =====================================================
\seccio{4}{Área del romboide y del triángulo}

\ex{1} \begin{tasks}(2)
  \task $12\cdot7=\res{84\text{ cm}^2}$
  \task $8,5\cdot4=\res{34\text{ m}^2}$
\end{tasks}

\ex{1} \begin{tasks}(3)
  \task $\frac{14\cdot9}{2}=\res{63\text{ cm}^2}$
  \task $\frac{7\cdot5}{2}=\res{17,5\text{ cm}^2}$
  \task $\frac{6,4\cdot3,5}{2}=\res{11,2\text{ m}^2}$
\end{tasks}

\ex{1} \begin{tasks}(1)
  \task $\frac{6\cdot4}{2}=\res{12\text{ cm}^2}$
  \task $\frac{8\cdot5}{2}=\res{20\text{ cm}^2}$
  \task $\frac{4\cdot5}{2}=\res{10\text{ cm}^2}$ (la altura cae fuera)
\end{tasks}

\ex{2} Los catetos hacen de base y altura: \ $A=\frac{9\cdot12}{2}=\res{54\text{ cm}^2}$. \ $P=9+12+15=\res{36\text{ cm}}$.

\ex{2} \begin{tasks}(1)
  \task $h=96:12=\res{8\text{ cm}}$
  \task $b\cdot9=90\Rightarrow b=\res{10\text{ m}}$
\end{tasks}

\ex{2} Base 8\,cm (sobre la línea de abajo) y altura 5\,cm: \ $A=\frac{8\cdot5}{2}=\res{20\text{ cm}^2}$.

\ex{3} $A=\frac{14\cdot12}{2}=\res{84\text{ cm}^2}$. \ Con base 15: $\frac{15\cdot h}{2}=84\Rightarrow 15\cdot h=168\Rightarrow h=\res{11,2\text{ cm}}$.

\ex{3} Todos tienen la misma base (10\,cm) y la misma altura (6\,cm, la distancia entre las paralelas): $A=\frac{10\cdot6}{2}=\res{30\text{ cm}^2}$ \rest{los tres}.


% =====================================================
\seccio{5}{Área del rombo y del trapecio}

\ex{1} \begin{tasks}(3)
  \task $\frac{16\cdot10}{2}=\res{80\text{ cm}^2}$
  \task $\frac{9\cdot6}{2}=\res{27\text{ m}^2}$
  \task $\frac{7,5\cdot4}{2}=\res{15\text{ cm}^2}$
\end{tasks}

\ex{1} \begin{tasks}(3)
  \task $\frac{(12+8)\cdot5}{2}=\res{50\text{ cm}^2}$
  \task $\frac{(9+4)\cdot6}{2}=\res{39\text{ m}^2}$
  \task $\frac{17\cdot4}{2}=\res{34\text{ cm}^2}$
\end{tasks}

\ex{2} $l=52:4=\res{13\text{ cm}}$; \ $A=\frac{24\cdot10}{2}=\res{120\text{ cm}^2}$

\ex{2} $A=\frac{(18+10)\cdot3}{2}=\res{42\text{ cm}^2}$; \ $P=18+10+5+5=\res{38\text{ cm}}$

\ex{2} \begin{tasks}(1)
  \task $15\cdot d=120\Rightarrow d=\res{8\text{ cm}}$
  \task $(14+10)\cdot h=144\Rightarrow 24\cdot h=144\Rightarrow h=\res{6\text{ m}}$
\end{tasks}

\ex{3} $A=\frac{1,2\cdot0,8}{2}=\res{0,48\text{ m}^2}=\res{4800\text{ cm}^2}$

\ex{3} $A=\frac{(85+55)\cdot40}{2}=2800\text{ m}^2=\res{0,28\text{ ha}}$. \ Precio: $2800\cdot12=\res{33\,600\text{ €}}$.


% =====================================================
\seccio{6}{Área de los polígonos regulares y de polígonos cualesquiera}

\ex{1} $P=5\cdot10=\res{50\text{ cm}}$; \ $A=\frac{50\cdot6,88}{2}=\res{172\text{ cm}^2}$

\ex{1} $P=6\cdot8=\res{48\text{ cm}}$; \ $A=\frac{48\cdot6,93}{2}=\res{166,32\text{ cm}^2}$

\ex{2} $l=40:8=\res{5\text{ cm}}$; \ $A=\frac{40\cdot6,04}{2}=\res{120,8\text{ cm}^2}$

\ex{2} $\frac{36\cdot a}{2}=93,6\Rightarrow 36\cdot a=187,2\Rightarrow a=\res{5,2\text{ cm}}$

\ex{2} Rectángulo $6\times4$ y triángulo de base 6 y altura 2: \ $24+\frac{6\cdot2}{2}=24+6=\res{30\text{ cm}^2}$

\ex{3} Rectángulo $7\times6=42$. Triángulos que sobran: $\frac{7\cdot2}{2}=7$, \ $\frac{4\cdot4}{2}=8$ \ y \ $\frac{3\cdot6}{2}=9$. \ $A=42-7-8-9=\res{18\text{ cm}^2}$

\ex{3} Baldosa: $A=\frac{60\cdot8,66}{2}=259,8\text{ cm}^2$. \ $25\text{ m}^2=250\,000\text{ cm}^2$; \ $250\,000:259,8\approx962,3$: hacen falta \rest{963 baldosas}.


% =====================================================
\seccio{7}{La circunferencia y su longitud}

\ex{1} \begin{tasks}(4)
  \task $a$: \rest{radio}
  \task $b$: \rest{diámetro}
  \task $c$: \rest{cuerda}
  \task $d$: \rest{arco}
\end{tasks}

\ex{1} \begin{tasks}(2)
  \task $2\cdot3,14\cdot4=\res{25,12\text{ cm}}$
  \task $2\cdot3,14\cdot9=\res{56,52\text{ m}}$
  \task $3,14\cdot20=\res{62,8\text{ cm}}$
  \task $3,14\cdot7=\res{21,98\text{ mm}}$
\end{tasks}

\ex{2} \begin{tasks}(1)
  \task $d=31,4:3,14=10$\,cm; \ $r=\res{5\text{ cm}}$
  \task $d=75,36:3,14=\res{24\text{ m}}$; \ $r=\res{12\text{ m}}$
\end{tasks}

\ex{2} $L=3,14\cdot70=\res{219,8\text{ cm}}$. \ $1\text{ km}=100\,000\text{ cm}$; \ $100\,000:219,8\approx455$ \ \rest{unas 455 vueltas}.

\ex{2} \begin{tasks}(3)
  \task $\frac{2\cdot3,14\cdot8\cdot90}{360}=\res{12,56\text{ cm}}$
  \task $\frac{2\cdot3,14\cdot12\cdot60}{360}=\res{12,56\text{ cm}}$
  \task $\frac{2\cdot3,14\cdot6\cdot150}{360}=\res{15,7\text{ cm}}$
\end{tasks}

\ex{3} \begin{tasks}(2)
  \task $\frac{3,14\cdot10}{2}+10=15,7+10=\res{25,7\text{ cm}}$
  \task $\frac{2\cdot3,14\cdot4}{4}+4+4=6,28+8=\res{14,28\text{ cm}}$
\end{tasks}

\ex{3} En 20 minutos gira $120^\circ$: \ $\frac{2\cdot3,14\cdot15\cdot120}{360}=\res{31,4\text{ cm}}$. \ En una hora: $2\cdot3,14\cdot15=\res{94,2\text{ cm}}$.


% =====================================================
\seccio{8}{El área del círculo y de las figuras circulares}

\ex{1} \begin{tasks}(2)
  \task $3,14\cdot36=\res{113,04\text{ cm}^2}$
  \task $3,14\cdot6,25=\res{19,625\text{ m}^2}$
  \task $3,14\cdot7^2=\res{153,86\text{ cm}^2}$
  \task $3,14\cdot15^2=\res{706,5\text{ cm}^2}$
\end{tasks}

\ex{2} \begin{tasks}(2)
  \task $\frac{3,14\cdot100}{2}=\res{157\text{ cm}^2}$
  \task $\frac{3,14\cdot36}{4}=\res{28,26\text{ cm}^2}$
\end{tasks}

\ex{2} \begin{tasks}(3)
  \task $\frac{3,14\cdot64\cdot45}{360}=\res{25,12\text{ cm}^2}$
  \task $\frac{3,14\cdot81\cdot120}{360}=\res{84,78\text{ cm}^2}$
  \task $\frac{3,14\cdot16\cdot270}{360}=\res{37,68\text{ cm}^2}$
\end{tasks}

\ex{2} \begin{tasks}(2)
  \task $3,14\cdot(49-16)=3,14\cdot33=\res{103,62\text{ cm}^2}$
  \task $3,14\cdot(100-36)=3,14\cdot64=\res{200,96\text{ m}^2}$
\end{tasks}

\ex{2} \begin{tasks}(1)
  \task $r^2=78,5:3,14=25\Rightarrow r=\res{5\text{ cm}}$
  \task $r=18,84:6,28=3$\,cm; \ $A=3,14\cdot9=\res{28,26\text{ cm}^2}$
\end{tasks}

\ex{3} Corona de radios $R=6$\,cm y $r=0,75$\,cm: \ $3,14\cdot(36-0,5625)\approx\res{111,27\text{ cm}^2}$

\ex{3} Grande: $3,14\cdot16^2=803,84\text{ cm}^2$. \ Dos pequeñas: $2\cdot3,14\cdot11^2=759,88\text{ cm}^2$. \ \rest{La pizza grande.}


% =====================================================
\seccio{9}{Figuras compuestas}

\ex{1} Completamos el rectángulo $10\times7$ y restamos el trozo $4\times3$: \ $A=70-12=\res{58\text{ cm}^2}$. \ $P=10+4+4+3+6+7=\res{34\text{ cm}}$.

\ex{1} $A=8\cdot5+\frac{8\cdot3}{2}=40+12=\res{52\text{ m}^2}$

\ex{2} Semicírculo de radio 2\,cm: \ $A=6\cdot4+\frac{3,14\cdot2^2}{2}=24+6,28=\res{30,28\text{ cm}^2}$. \ $P=6+6+4+\frac{2\cdot3,14\cdot2}{2}=16+6,28=\res{22,28\text{ cm}}$

\ex{2} Los cuatro cuartos forman un círculo entero de radio 4\,cm: \ $A=8^2-3,14\cdot4^2=64-50,24=\res{13,76\text{ cm}^2}$

\ex{2} Rectángulo exterior: $(10+3)\times(6+3)=13\cdot9=117\text{ m}^2$. \ $A=117-60=\res{57\text{ m}^2}$

\ex{2} Los dos semicírculos forman un círculo de radio 25\,m: \ $A=80\cdot50+3,14\cdot25^2=4000+1962,5=\res{5962,5\text{ m}^2}$. \ $P=2\cdot80+3,14\cdot50=160+157=\res{317\text{ m}}$

\ex{3} Centro: $3,14\cdot9=\res{28,26\text{ cm}^2}$. \ Zona del medio: $3,14\cdot(36-9)=\res{84,78\text{ cm}^2}$. \ Zona exterior: $3,14\cdot(81-36)=\res{141,3\text{ cm}^2}$.

\ex{3} La diagonal divide la hoja en dos mitades. Cada mitad es un cuarto de círculo menos un triángulo: $\frac{3,14\cdot36}{4}-\frac{6\cdot6}{2}=28,26-18=10,26$. \ $A=2\cdot10,26=\res{20,52\text{ cm}^2}$

\ex{3} $A=1,2\cdot1,5+\frac{3,14\cdot0,6^2}{2}=1,8+0,5652=2,3652\text{ m}^2$. \ $2,3652\cdot45\approx\res{106,43\text{ €}}$


% =====================================================
\seccio{10}{El teorema de Pitágoras}

\ex{1} \begin{tasks}(1)
  \task Catetos $x$ e $y$; hipotenusa $z$: \ $\res{z^2=x^2+y^2}$
  \task Catetos $n$ y $p$; hipotenusa $m$: \ $\res{m^2=n^2+p^2}$
  \task Catetos $r$ y $s$; hipotenusa $t$: \ $\res{t^2=r^2+s^2}$
\end{tasks}

\ex{1} \begin{tasks}(2)
  \task $\sqrt{9+16}=\sqrt{25}=\res{5\text{ cm}}$
  \task $\sqrt{36+64}=\sqrt{100}=\res{10\text{ m}}$
  \task $\sqrt{25+144}=\sqrt{169}=\res{13\text{ cm}}$
  \task $\sqrt{81+144}=\sqrt{225}=\res{15\text{ dm}}$
\end{tasks}

\ex{1} \begin{tasks}(2)
  \task $\sqrt{100-36}=\sqrt{64}=\res{8\text{ cm}}$
  \task $\sqrt{169-144}=\sqrt{25}=\res{5\text{ m}}$
  \task $\sqrt{289-225}=\sqrt{64}=\res{8\text{ cm}}$
  \task $\sqrt{625-49}=\sqrt{576}=\res{24\text{ cm}}$
\end{tasks}

\ex{2} \begin{tasks}(2)
  \task $c=\sqrt{4+25}=\sqrt{29}\approx\res{5,39\text{ cm}}$
  \task $c=\sqrt{16+36}=\sqrt{52}\approx\res{7,21\text{ cm}}$
  \task $b=\sqrt{64-9}=\sqrt{55}\approx\res{7,42\text{ cm}}$
  \task $b=\sqrt{100-49}=\sqrt{51}\approx\res{7,14\text{ cm}}$
\end{tasks}

\ex{2} \begin{tasks}(2)
  \task $10^2=100$; \ $36+64=100$: \ \rest{sí}
  \task $9^2=81$; \ $25+49=74$: \ \rest{no}
  \task $17^2=289$; \ $64+225=289$: \ \rest{sí}
  \task $15^2=225$; \ $100+121=221$: \ \rest{no}
\end{tasks}

\ex{2} La escalera es la hipotenusa: \ $h=\sqrt{5^2-1,4^2}=\sqrt{25-1,96}=\sqrt{23,04}=\res{4,8\text{ m}}$

\ex{2} Los dos trayectos son los catetos: \ $d=\sqrt{12^2+9^2}=\sqrt{144+81}=\sqrt{225}=\res{15\text{ km}}$

\ex{3} Triángulo rectángulo de catetos 6\,cm y 8\,cm: \ $AB=\sqrt{36+64}=\sqrt{100}=\res{10\text{ cm}}$

\ex{3} Cada cable: $\sqrt{12^2+5^2}=\sqrt{169}=13$\,m. \ Cuatro cables: $4\cdot13=\res{52\text{ m}}$


% =====================================================
\seccio{11}{Pitágoras para calcular áreas y perímetros}

\ex{1} \begin{tasks}(3)
  \task $\sqrt{64+36}=\res{10\text{ cm}}$
  \task $\sqrt{576+49}=\res{25\text{ m}}$
  \task $\sqrt{16+16}=\sqrt{32}\approx\res{5,66\text{ cm}}$
\end{tasks}

\ex{1} La altura divide la base en dos mitades de 8\,cm: \ $h=\sqrt{10^2-8^2}=\sqrt{36}=\res{6\text{ cm}}$. \ $A=\frac{16\cdot6}{2}=\res{48\text{ cm}^2}$. \ $P=16+10+10=\res{36\text{ cm}}$

\ex{2} Las semidiagonales miden 9 y 12\,cm: \ $l=\sqrt{81+144}=\sqrt{225}=\res{15\text{ cm}}$. \ $P=4\cdot15=\res{60\text{ cm}}$. \ $A=\frac{18\cdot24}{2}=\res{216\text{ cm}^2}$

\ex{2} Triángulo de catetos $12-7=5$ y 12: \ $c=\sqrt{25+144}=\sqrt{169}=13$\,cm. \ $P=12+7+12+13=\res{44\text{ cm}}$. \ $A=\frac{(12+7)\cdot12}{2}=\res{114\text{ cm}^2}$

\ex{2} A cada lado sobran $(16-4):2=6$\,cm: \ $h=\sqrt{100-36}=\res{8\text{ cm}}$. \ $A=\frac{(16+4)\cdot8}{2}=\res{80\text{ cm}^2}$. \ $P=16+4+10+10=\res{40\text{ cm}}$

\ex{2} El radio es igual al lado (4\,cm): \ $a=\sqrt{4^2-2^2}=\sqrt{12}\approx\res{3,46\text{ cm}}$. \ $P=6\cdot4=24$\,cm; \ $A=\frac{24\cdot3,46}{2}=\res{41,52\text{ cm}^2}$

\ex{3} Lado: $18:3=6$\,cm. \ $h=\sqrt{6^2-3^2}=\sqrt{27}\approx\res{5,20\text{ cm}}$. \ $A=\frac{6\cdot5,2}{2}=\res{15,6\text{ cm}^2}$

\ex{3} La diagonal del cuadrado es un diámetro (8\,cm): \ $l^2+l^2=64\Rightarrow l^2=32$, que es el área del cuadrado. \ $A=3,14\cdot4^2-32=50,24-32=\res{18,24\text{ cm}^2}$

\ex{3} Cada lado del tejado: $\sqrt{4^2+3^2}=\sqrt{25}=5$\,m. \ $P=8+5+5+5+5=\res{28\text{ m}}$. \ $A=8\cdot5+\frac{8\cdot3}{2}=\res{52\text{ m}^2}$


\end{document}
